\quellseite{688}
\part{III. Auslegungen über das vierte Buch Mosis}

\chapter{Der Segen}

\begin{center}
so man nach der Messe spricht über das Volk, aus dem vierten Buch Mosis am sechsten Capitel, ausgelegt durch D. Martin Luther.\footnote{Luther predigte in den zwanziger Jahren fortlaufend über alle fünf Bücher Mosis, wie wir aus dem Anfang der Schrift: „Auslegung über etliche Capitel des fünften Buchs Mose“ ersehen. Doch von den Predigten über das dritte Buch Mose ist uns nichts erhalten, und von dem vierten Buche nur die gegenwärtige Schrift. Sie erschien im Jahre 1532 unter dem von uns darüber gesetzten Titel in zwei Einzelausgaben; eine bei Nickel Schirlentz in Wittenberg, die andere bei Kunigund Hergotin in Nürnberg. Falls diese Schrift zu den fortlaufenden Predigten Luthers gehört, können wir mit ziemlicher Sicherheit annehmen, daß sie im Jahre 1527 gepredigt sein wird. Bei dieser Zeitbestimmung gründen wir uns auf folgende: In Buchwalds Andreas Poachs handschriftliche Sammlung ungedruckter Predigten D. M. Luthers finden wir Bd. I erste Hälfte, S. XXIX in dem Index zu Poachs verlorengegangenen Predigten im Jahre 1528 zwanzig Predigten verzeichnet, welche Luther über das vierte Buch Mosis gehalten hat: am 8. März 1528 predigte er über das 15. Capitel dieses Buches, und am 13. December 1528 stand er bei dem 31. Capitel. Daher mögen wir schließen, daß Luther im Jahre 1527 über das sechste Capitel gepredigt habe. Weil Luther nur Nachmittags über die Bücher Mosis predigte, der Index aber in den dem Jahre 1528 vorhergehenden Jahren nur selten eine Nachmittags-predigt aufführt, so fehlt in demselben die Nachricht darüber. In den Sammlungen findet sich diese Schrift: Wittenberger (1556), Bd. V, Bl. 245; Jenaer (1566), Bd. V, Bl. 480 b; Altenburger, Bd. V, S. 956; Leipziger, Bd. IV, S. 89 und Erlanger, Bd. 36, S. 155. Diese Schrift ist auch ins Lateinische übersetzt und findet sich in der lateinischen Wittenberger Ausgabe (1554), Tom. III, fol. 11 b. Wir geben den Text nach der Erlanger Ausgabe, die den Urtext bringt, unter Vergleichung der Wittenberger und der Jenaer.}

Gepredigt wahrscheinlich 1527, herausgegeben 1532.
\end{center}

\begin{schriftzitat}
4 Mos. 6, 22–27. Und der HErr sprach zu Mose: Sage Aaron und seinen Söhnen: Also sollt ihr die Kinder Israel segnen, und sprechen: Der HErr segne dich, und behüte dich. Der HErr erleuchte sein Angesicht über dir, und sei dir gnädig. Der HErr erhebe sein Angesicht über dich, und gebe dir Friede. Denn sie sollen meinen Namen legen auf die Kinder Israel, und ich will sie segnen. Der HErr segne dich, und behüte dich.
\end{schriftzitat}

1. Das erste Stück dieses Segens betrifft das leibliche Leben und Güter, wie 1 Mos. 1, 28. und auch sonst oft gesagt wird: „Und GOtt segnete sie, und sprach: Wachset, und seid fruchtbar.“ Daher auch in der Schrift (1 Mos. 27, 34. 36. 38.) die gemeine Weise ist zu reden: „Gib mir einen Segen.“ „Hast du nicht mehr Segen?“ das ist: Gib mir etwas, als Gut, Brod, Kleid \&c. Denn es ist alles eitel GOttes Gabe, und durch seinen Segen haben wir, was wir haben; und heißt auch darum ein Segen, das ist, eine GOttes Gabe, die er uns durch seinen Segen gibt.

2. So wünschet nun dieser Segen dem Volke, daß ihm GOtt wolle Glück und alles Gut geben, und auch behüten und erhalten, nämlich, daß es erstlich an ihm selbst leiblich fruchtbar sei, treu ehelich Gemahl, Nahrung, Kleider, und alles, was noth ist zu diesem leiblichen Leben, habe, es sei Haus, Hof, Acker, Vieh, Gesinde. Und wenn er es nun gibt, und wir es haben, daß er es auch behüte und bewahre; als, den Leib vor Krankheit und Plagen, das Vieh, Haus, Acker vor Feuer, Wasser, Ungewitter und allerlei Schaden.

3. Und lehrt uns solcher Segen daneben, daß wir dankbar sein sollen, und erkennen, daß unser Leib und Leben, sammt allen Gütern und Nothdurft, nicht durch unser Werk noch Arbeit gemacht noch gewonnen werden mögen, auch nicht durch unsern Fleiß und Sorge bewahrt noch erhalten werden, sondern es ist alles GOttes Segen, Gottes Hül und Sorge [Ps. 127, 1. 2.].
\quellseite{689}
nicht, daß wir darum sollten müßig gehen, nichts arbeiten, noch hüten, sondern sollen das Unsere thun, mit Arbeiten und Hüten, aber doch wissen, daß GOtt müsse das Gedeihen und Gerathen geben in unserer Arbeit und Hut, wie St. Paulus lehrt 1 Cor. 3, 6.: „Ich habe gepflanzt, Apollo hat begossen; aber GOtt hat das Gedeihen gegeben.“

4. Denn wir sehen, wie unzählige Krankheit und Gebrechen unserm eigenen Leibe, Weibern, Kindern, Gesinde, wie unzählige Jahr und Plage dem Vieh, wie unzählige Schade und Ungeziefer den Früchten und andern Gütern widerfahren kann, und täglich widerfährt, daß nicht möglich ist, aus unserer Macht oder Witz ein Hühnlein oder Ferklein, ja auch nicht ein Körnlein oder Hälmlein zu erhalten oder behüten, geschweige, daß wir es machen oder schaffen sollten. Der Schöpfer, der es uns alles schafft und gibt, der muß auch alles behüten und erhalten; ohne daß wir arbeiten müssen, und nicht müßig noch laß gehen.

\begin{schriftzitat}
V. 25. Der HErr erleuchte sein Angesicht über dir, und sei dir gnädig.
\end{schriftzitat}

5. Das andere Stück betrifft das geistliche Wesen und die Seele, denn solches Segens ist das Vieh und Güter nicht empfänglich, sondern bleiben in dem ersten Stück, da sie mit uns nach dem Leibe den gemeinen Segen haben.

6. Und hier müssen wir Hebräisch lernen, und der Sprache gewohnen. Denn es lautet im Deutschen gar nichts, wenn ich spreche: „GOtt erleuchte sein Angesicht über dir.“ Und ist doch nicht wohl anders zu geben noch zu verdeutschen, und müssen es also hebräisch nachreden, und bleiben lassen. Und ist auf deutsch so viel gesagt: GOtt der HErr erzeige sich dir freundlich und tröstlich, sehe dich nicht sauer an noch zornig, erschrecke dein Herz nicht, sondern lache dich fröhlich und väterlich an, daß du fröhlich und getrost von ihm werdest, und eine freudige, herzliche Zuversicht zu ihm habest. Welches geschieht, so er unsere Sünde uns vergibt, und nicht mit uns rechnet, sondern durch sein Wort und Geist unser betrübtes, blödes Gewissen los spricht, daß unser Herz fühlen könne, wie GOtt unsere Sünde nicht ansehen, sondern vergessen will ewiglich, und hinfort immerdar gnädig, freundlich, tröstlich und sanft gegen uns sein.

7. Gleichwie die liebe Sonne, wenn sie aufgeht und streuet ihren reichen Schein und mildes Licht in alle Welt, so thut sie nicht anders, denn daß sie erleuchtet ihr Angesicht über alle Welt; das ist, sie scheint helle und fröhlich, und erfreuet nicht allein die Menschen, sondern Thiere, Vögel etc., und alles, was im Himmel und Erden ist, macht sie licht und gleich als neu geschaffen. Wenn sie aber ihr klares, liebliches Angesicht verbirgt oder untergeht, so sieht alles Ding sauer, betrübt, finster, schwarz, als wäre es todt und in der Hölle.

8. Also sollst du hier den Text des Segens auch verstehen. Wenn GOtt sein Wort gibt, so läßt er sein Angesicht fröhlich und helle scheinen über alle Gewissen, und macht sie damit fröhlich, fest, licht, und als gar neue Herzen und neue Menschen. Denn es bringt Vergebung der Sünden, und zeigt GOtt als einen gnädigen, barmherzigen Vater an, welchen unser Leid und Betrübniß jammert und erbarmt. Wenn er aber sein Wort wegnimmt, da verbirgt er sein Angesicht, da wird es denn finster, schwarz, betrübt, Jammer und Leid im Gewissen, da fühlt man eitel Zorn und keine Gnade.

9. So wünscht nun dieser Segen (wenn ich es deutsch aussprechen soll), daß uns GOtt wolle sein gnädiges Wort reichlich lassen scheinen und damit fröhliche Gewissen machen, daß uns die Sünde, Gewissen, Gesetz und Teufel nicht erschrecken, noch verzagt, blöde und betrübt machen, daß wir nicht seinen Zorn, sondern seine Gunst und Liebe in unserm Herzen fühlen. Das ist der rechte, geistliche, große Segen, der ein geistliches und ewiges Leben gibt, welcher den Thieren, Vieh, Vögeln etc. nicht kann gegeben werden.

10. Er setzt aber dazu: „Und sei dir gnädig“, und läßt es nicht genug sein, daß er sein Angesicht über uns scheinen oder leuchten läßt. Denn es sind auch zwei Stücke in diesem Segen, welche St. Paulus pflegt unterschiedlich zu nennen, gratia et donum, Gnade und Gabe. Das erste ist Gnade oder Gunst, wenn er seinen Zorn wendet und uns durch sein Wort gnädiglich ansieht, und uns alle Sünde vergibt, und also unser Herz gegen ihm sicher und fröhlich macht. Darnach gibt er auch die reichen und mancherlei Gaben seines Geistes, dadurch wir viel ausrichten, beide in uns selbst und in aller Welt; als, lehren, trösten, rathen, helfen, taufen, Sacrament reichen, Wunder thun, und dergleichen; wie 1 Cor. 12, 4. f. Paulus erzählt.

\quellseite{690}
11. Gleichwie die liebe Sonne, wenn sie aufgeht, macht sie nicht allein fröhlich alle Welt mit ihrem lieblichen, schönen Lichte, sondern wirkt, gibt und hilft, daß beide, Menschen und Vieh, allerlei thun, arbeiten, gehen, stehen, und alle Glieder brauchen können, erwärmt und treibt fort allerlei Gewächse und Früchte, bis sie alles reif und vollkommen macht. Also spricht hier der Segen auch, daß GOtt uns nicht allein wollt helle scheinen mit dem Worte der Gnaden, und uns fröhlich machen durch Vergebung der Sünden und Anzeigung seiner Gunst gegen uns, sondern auch uns gnädig sein, das ist, gnädiglich nachdrücke, mit Geist und Gaben stärke, zu thun allen seinen Willen. Darum heißt dies Wort „gnädig sein“ so viel als begnaden oder gnädiglich geben; wie der Erzvater Jakob, 1 Mos. 33, 5., solches Wortes braucht, und spricht zu Esau: Dies sind meine Kindlein, die mir der HErr begnadet, beschert, oder gnädiglich gegeben hat.

12. Und lehrt uns hier abermal dieser Segen, daß wir dankbar sein sollen, und erkennen, wie wir unsere Sünde durch kein Verdienst noch Werke können los werden, auch alle Heiligkeit und Weisheit auf Erden, alle anderen Prediger und Lehrer, sie seien wer sie wollen, nicht ein einig Gewissen in der allergeringsten Sünde stillen noch trösten können, sondern es muß das Licht göttliches Worts allein thun. Nicht, daß wir sollten darum keine guten Werke thun, oder verrucht dahin leben, sondern wir sollen uns in guten Werken üben, und doch wissen, daß ein fröhlich Gewissen keines andern Weges herkomme, denn wo GOtt sein Angesicht über uns leuchten läßt, das ist, durch Vergebung der Sünde, aus lauter Gnade und Barmherzigkeit; dazu auch weder predigen, oder bekennen, noch ein einiges recht christliches Werk oder Wort anzuführen vermögen, geschweige denn zu vollenden, wo er uns nicht gnädig ist, das ist, wo er uns nicht gnädiglich mit seinen Gaben begnadet, stärkt und treibt.

\begin{schriftzitat}
V. 26. Der HErr erhebe sein Angesicht über dich, und gebe dir Friede.
\end{schriftzitat}

13. Das dritte Stück betrifft auch das geistliche Wesen und die Seele, und ist ein Wunsch des Trostes und endlichen Sieges unter dem Kreuz, Tod, Teufel und aller höllischen Pforten, sammt der Welt und übrigen bösen Lüsten unseres Fleisches. Denn ob er uns wohl gnädig ist worden, die Sünde vergeben hat und mit seinem Geiste geziert, so haben wir gleichwohl noch vor uns, ja, liegen noch stets im Kampfe wider den Teufel und übrige Sünden.

14. Denn auch der Teufel, wenn er gewahr wird, daß wir GOttes gnädiges Wort haben, und aus seinem Reiche gerissen sind, und GOttes Kinder worden, so wird er auch toll und thöricht, und hebt an uns der Welt Gewalt mit Schwert, Feuer etc., dazu der Welt Weisheit mit Ketzerei und Irrthum mannigfältiglich; auch der Welt Heiligkeit durch falsche Brüder und Heuchler: daß wohl noth ist, nicht allein ansehen uns zu erleuchten, und [den] Geist [zu] geben, oder wie St. Paulus Röm. 8, 23. sagt, die Erstlinge des Geistes, sondern halte an über uns, daß wir die Fülle des Geistes, und endlich den ganzen Sieg erlangen. Denn nicht, wer da anfähet, sondern, wer da beharret, wird selig werden, Matth. 24, 13.

15. Und hier müssen wir abermal hebräisch lernen reden, denn es lautet auf deutsch nichts, wenn ich spreche: „Der HErr erhebe sein Angesicht über dich.“ Und man kann es doch nicht wohl deutsch geben. Denn die Meinung soll die sein: Wenn GOtt sein Angesicht über uns leuchten oder scheinen läßt, so machen doch die Anfechtung und Verfolgung des Teufels, der Welt und unseres Fleisches ein solch Ungewitter in unsern Herzen, daß uns dünkt, das liebe Licht des theuern göttlichen Worts wolle untergehen, und uns im Finstern lassen.

16. Gleich als wenn die Sonne wohl hübsch und fein aufgeht, fröhlich scheint, und lieblich leuchtet, und doch ein Gewölke und Ungewitter sich wider sie erhebt und ihr den Schein nimmt, daß man keine Sonne mehr sieht, und ist, als wolle sie versinken, und sich das Ungewitter überwinden und unterdrücken lassen, daß man wohl mag sagen zur Sonne: Halte fest, liebe Sonne, und laß dich die Wolken und Wetter nicht unterdrücken, oder den Tag zur Nacht machen; sondern erhebe dein schönes Licht über alle Wolken und Wetter, und erhalte uns den Tag, daß nicht die Wolken und Wetter mit ihrem Finsterniß den Sieg behalten, sondern du mit deinem schönen Lichte obliegest und Oberhand behaltest. Also hier auch will dieser Segen wünschen, daß GOtt der HErr das Licht seines Worts wollt erheben über uns, und also darüber halten, daß
\quellseite{691}
es höher und stärker leuchte in unserm Herzen, denn alle Anfechtung des Teufels, Todes und der Sünde, Verzweifeln, Verzagen, Erschrecken und alles Unglück sein kann. Denn wo er das nicht thut, so ist uns der Teufel mit seinem Wetter und Wolken zu mächtig, und verfinstert und verdunkelt uns das liebe angefangene Licht seines Worts, und bringt uns so jämmerlich darum, daß [es] hernach ärger mit uns wird, denn es vor je gewesen ist; wie wir sehen und erfahren täglich an denen, die vom Worte fallen, durch Gewalt gedrungen oder durch Ketzerei verführt.

17. Denn, wie Christus [Luc. 11, 26.] sagt: Wenn der unreine Geist ausgetrieben ist, so kommt er mit sieben ärgeren wieder, daß wohl noch thut hier, segnen und beten, rufen und wünschen, GOtt wolle sein Angesicht über uns erheben, das ist, das Licht seines werthen Worts in unserm Herzen gewaltiglich erhalten und oben schweben lassen wider den Mörder und Lügner, den Teufel, welcher durch Mord und Lügenlehre dasselbe Wort in uns unterdrücken und dämpfen will.

18. Und dies dritte Stück dieses Segens begreift auch zwei, nämlich, Erheben des Angesichts, und den Frieden. Denn er wünscht nicht allein, daß wir gestärkt und getröstet werden durch Erhebung seines Angesichts, wider den Teufel etc., sondern, daß er uns auch ein friedlich Herz und guten Muth verleihe in solchem Kampfe und Unruhe, auf daß wir nicht allein erdulden und endlich obliegen, sondern auch mitten in dem Kampfe und Unruhe Frieden haben mögen, GOtt loben und danken, und nicht murren noch ungeduldig werden wider seinen göttlichen Willen, wie St. Paulus [Col. 3, 15.] lehrt: daß der Friede in unserm Herzen solle den Sieg behalten, daß wir nichts wider GOtt noch Menschen durch Ungeduld vornehmen, sondern beide, inwendig und auswendig, gegen GOtt und den Menschen, still und friedlich bleiben, bis der endliche und ewige Friede komme.

19. Dieser Segen ist nicht weit von dem andern gemeinen Segen, den man in der lateinischen Sprache gibt, und auf deutsch also lautet: Es segne euch GOtt der Vater, und der Sohn, und Heiliger Geist, Amen. Denn dem Vater wird zugeeignet das Werk der Schöpfung, welches dieser unser Segen auch rührt, und klarer ausdrückt, da er spricht: „Der HErr segne dich, und behüte dich“, das ist, er gebe dir gnädiglich Leib und Leben, und was dazu gehört. Also, dem Sohne wird zugeeignet das Werk der Erlösung, welches dieser Segen auch rührt und erklärt, da er spricht: „Der HErr erleuchte sein Angesicht über dir“ etc., das ist, er helfe dir von Sünden, und sei dir gnädig, und gebe dir seinen Geist. Und dem Heiligen Geist wird zugeeignet das Werk der täglichen Heiligung, Trost und Stärke wider den Teufel, und endlich die Auferweckung vom Tode, welches dieser Segen auch rührt und erklärt, da er spricht: „Der HErr erhebe sein Angesicht“ etc., das ist, er wolle dich stärken, trösten, und endlich den Sieg geben, wie droben [§ 13 ff.] gesagt.

20. Daß aber GOtt gebeut bei diesem Segen, sie sollen seinen Namen über die Kinder Israel brauchen, wenn sie segnen, hat die Meinung, daß er verboten haben will die Segen der Abgötter, als da war der Gott Baal, der Gott zu Dan, und dergleichen, mit welchen die falschen Propheten, unter dem Schein des rechten GOttes, das Volk pflegten zu segnen, wie man hin und wieder in den Propheten klagen hört; gleichwie bei uns die Pfaffen und Mönche uns gesegnet haben durch der Heiligen Namen, wie man findet in den Collecten, da sie unverschämt beten: GOtt wolle uns um dieses und des Heiligen willen dies und das geben. Das sind abgöttische Gebete und Segen, die nicht im Namen GOttes allein geschehen, wie es doch sein soll etc.
